\section{Finite-width realization of a planted scalar channel}
\label{sec:critical-finite-width-channel}

The construction below gives a finite-width lower bound on the raw
probability of a target.  It follows a scalar component and diagonal
variances, rather than approximating an entire hidden covariance matrix.
In particular its width depends on the logarithm of the input-set size.

Let $X=\{-1,1\}^d$, $m=|X|$, $H=L-1\ge1$, and let
$f:X\to\{-1,1\}$.  We first specify the controlled data used by the
construction.  All auxiliary variables below are absent from the final
network; they are used only in its probabilistic construction.

\begin{definition}[A scalar-channel program]
\label{def:critical-scalar-channel-program}
The initial controlled row law, including auxiliary variables, has entropy
at most $D_1$ relative to a reference law whose weight marginal is the
original first-row Gaussian.  Write $h(x)=\tanh(w\cdot\bar x)$.
It has a score $b_1$ and numbers $v_1>0,r_1\ge0$ satisfying
\[
 \mathbb E b_1^2=1,\quad |b_1|\le B,\quad
 \mathbb E[b_1h(x)]=\sqrt{v_1}f(x),\quad
 \mathbb E h(x)^2=v_1+r_1\le1.
\]
For $2\le j\le H$, a controlled law of $(Z,U)$ has entropy at
most $D_j$ relative to a reference law whose $Z$ marginal is
$\mathcal N(0,1)$.  With an independent standard Gaussian $G$, put
\[
 T_{j,x}=\tanh\!\left(\sqrt{v_{j-1}}f(x)Z+
                                      \sqrt{r_{j-1}}G\right).
\]
There is a score $b_j=b_j(Z,U)$ such that
\[
 \mathbb E b_j^2=1,\quad |b_j|\le B,\quad
 \mathbb E[b_jT_{j,x}]=\sqrt{v_j}f(x),\quad
 \mathbb E T_{j,x}^2=v_j+r_j\le1,
\]
where $v_j>0$ and $r_j\ge0$.  Define the cost and accumulated
entropy distortion by
\[
 A=\sum_{j=1}^H D_j+\frac1{8v_H},\qquad
 S=\sum_{j=2}^H\sqrt{D_j},\qquad B\ge1.
\]
The auxiliary part of each reference law may depend on its Gaussian
coordinate.  Its marginal on the actual incoming weights is always the
original Gaussian law.
\end{definition}

\begin{theorem}[A quantitative raw target-probability bound]
\label{thm:critical-finite-channel-mass}
There is a universal $C$ such that a scalar-channel program has the
following consequence.  For every $\tau>0$ and every integer width
\begin{equation}\label{eq:critical-channel-width}
 N\ge C B^4H^4 e^{4S}(A+1)^6(1+\tau^{-1})^4
       \left[\log(2m)+\log(H+A+\tau^{-1}+3)\right],
\end{equation}
the original critical Gaussian network satisfies
\[
 \mu_N\{f(x)z_w(x)>1/2\text{ for all }x\in X\}
       \ge e^{-N(A+\tau)}.
\]
No conditioning or support restriction is added to this prior.
\end{theorem}

\begin{proof}
We give the metric, the finite realization, and the probability comparison
separately.  All expectations in the construction refer to the indicated
controlled law.

\paragraph{Projection onto a bounded score.}
For bounded features $h(x)$ and a scalar score $b$ with
$\mathbb E b^2=1$, set
\[
 c(x)=\mathbb E[bh(x)],\qquad
 q(x)=\mathbb E h(x)^2,\qquad
 s(x)=\sqrt{q(x)-c(x)^2}.
\]
For $N$ independent rows let
\[
 M=\frac1N\sum_i b_i^2,\quad
 \widehat c(x)=\frac{N^{-1}\sum_i b_i h_i(x)}{\sqrt M},\quad
 \widehat q(x)=\frac1N\sum_i h_i(x)^2,
 \quad \widehat s(x)=\sqrt{\widehat q(x)-\widehat c(x)^2}.
\]
Suppose $|M-1|$, all errors in $N^{-1}\sum_i b_i h_i(x)$,
and all errors in $\widehat q(x)$ are at most $\alpha\le1/4$.
Then Cauchy--Schwarz and elementary square-root estimates give
\begin{equation}\label{eq:critical-channel-sampling-metric}
 \max_x\sqrt{(\widehat c(x)-c(x))^2+
                         (\widehat s(x)-s(x))^2}
       \le8\sqrt\alpha.
\end{equation}
Indeed $|\widehat c-c|\le4\alpha$, $|c|\le1$, and
$|(\widehat q-\widehat c^2)-(q-c^2)|\le13\alpha$.
The inequality $(\sqrt u-\sqrt v)^2\le|u-v|$ finishes the bound.

The coefficients $a_i=b_i/(N\sqrt M)$ satisfy
$\|a\|_2^2=1/N$.  Consequently the actual empirical feature kernel
has the exact decomposition
\begin{equation}\label{eq:critical-channel-exact-component}
 K(x,y)=\frac1N\sum_i h_i(x)h_i(y)
       =\widehat c(x)\widehat c(y)+R(x,y),\qquad R\succeq0.
\end{equation}
No approximate covariance inequality is needed.

\paragraph{Propagation of the component metric.}
Suppose the current finite features have coefficient $c(x)=a^Th(x)$,
$\|a\|^2=1/N$, and residual diagonal $s(x)^2$ in
\eqref{eq:critical-channel-exact-component}.  Compare them with the
ideal coefficient $\sqrt v f(x)$ and residual diagonal $r$ using
\[
 e=\max_x\sqrt{(c(x)-\sqrt v f(x))^2+(s(x)-\sqrt r)^2}.
\]
Set $u=\sqrt N a$, a unit vector.  Sample an incoming Gaussian
coordinate $Z$ with the prescribed controlled auxiliary law, and an
independent standard Gaussian vector $\xi\in\mathbb R^N$.  The
actual incoming row is
\[
 w=\{Zu+(I-uu^T)\xi\}/\sqrt N.
\]
Under its reference law this row is exactly $\mathcal N(0,I/N)$.
Its preactivation at $x$ is distributed as $c(x)Z+s(x)G$, with $G$
standard Gaussian independent of $(Z,U)$.  For comparison with the
ideal preactivation use the same $Z,U,G$ at this one input.

Data processing and Gaussian entropy duality imply
\[
 \max\{1,\mathbb E Z^2\}\le e^{2\sqrt{D_j}}.
\]
For completeness, writing $t=\mathbb E Z^2$, one has
$D_j\ge(t-1-\log t)/2$; for $t\ge1$, the latter is at least
$(\log t)^2/4$.  Independence and the zero mean of $G$ therefore
bound the preactivation RMS difference by $e^{\sqrt{D_j}}e$.
The same bound holds after tanh.

Projection onto the common unit vector $b_j$ in the coupled
$L^2$ space is contractive.  The difference of the norms of its two
orthogonal residuals is bounded by their $L^2$ distance.  Hence the
population next coefficient and residual RMS, compared with
$\sqrt{v_j}f$ and $\sqrt{r_j}$, have metric error at most
$e^{\sqrt{D_j}}e$.  This argument uses only one-input marginal
couplings; the actual sampled network has its usual joint law on all
inputs.  No simultaneous choice of the comparison residual correlations
is required.

Applying \eqref{eq:critical-channel-sampling-metric} at the new
layer now gives the recursion
\begin{equation}\label{eq:critical-channel-metric-recursion}
 e_1\le8\sqrt\alpha,\qquad
 e_j\le e^{\sqrt{D_j}}e_{j-1}+8\sqrt\alpha,
 \qquad e_H\le8H e^S\sqrt\alpha.
\end{equation}

\paragraph{Uniform concentration at finite width.}
At every stage, conditional on the complete previous construction,
the new rows and their scores are independent.  The observables
$b^2$, $bh(x)$, and $h(x)^2$ are bounded by $B^2,B,1$ in absolute
value.  Hoeffding's inequality and a union bound over the $2m+1$
observables, then over the $H$ stages, show that all the required
moment estimates hold with probability at least $1-\kappa/2$ if
\begin{equation}\label{eq:critical-channel-hoeffding}
 N\ge C B^4\alpha^{-2}\log\frac{8H(2m+1)}\kappa.
\end{equation}
The same assertion follows by successive conditioning, since the
bound is uniform in the realized previous layers.  Stop using the
controlled construction if a score Gram $M$ is outside $[1/2,3/2]$;
thereafter use the original prior.  This defines the construction on
every outcome, and the good moment event never triggers stopping.

Choose
\[
 \kappa=\min\left\{\frac14,\frac\tau{32(A+1)}\right\},\qquad
 \alpha=\left(\frac{\kappa\sqrt{v_H}}{64He^S}\right)^2.
\]
On the good event, \eqref{eq:critical-channel-metric-recursion} gives
$e_H\le\kappa\sqrt{v_H}/8$.  In particular the final coefficient
satisfies $f(x)c(x)\ge\sqrt{v_H}(1-\kappa/8)$ everywhere.

\paragraph{Readout and change of measure.}
When this final inequality holds, shift the Gaussian readout mean to
$ta$, where $\|a\|^2=1/N$ is its final score coefficient vector and
\[
 t=\frac{1+\kappa}{2\sqrt{v_H}}.
\]
The mean output then has signed margin at least $1/2+\kappa/4$.
At all other histories use zero readout mean.  The conditional readout
noise has variance $K(x,x)/N\le1/N$ at every input.  A Gaussian tail
bound and union bound make its maximum absolute value smaller than
$\kappa/8$ with probability at least $1-\kappa/2$, provided
$N\ge C\kappa^{-2}\log(4m/\kappa)$.  The complete controlled
construction, denoted $Q_N$, therefore produces the strict target
event $E_f$ with probability at least $1-\kappa$.

Augment the original network prior with the same reference auxiliary
kernels used above.  Conditional on the complete past, each reference
incoming weight row remains its original Gaussian, including under
the history-dependent orthogonal decomposition.  Chain rule and data
processing bound each controlled hidden block's relative entropy by
$ND_j$; stopping can only reduce this bound.  The shifted readout
has relative entropy $Nt^2/2$.  Consequently, after discarding all
auxiliary variables,
\[
 D(Q_N\Vert\mu_N)\le
 N\left[\sum_jD_j+\frac{(1+\kappa)^2}{8v_H}\right]
 \le N(A+3A\kappa).
\]
Binary relative entropy for the event $E_f$ gives
\[
 -\frac1N\log\mu_N(E_f)
 \le\frac{A+3A\kappa+N^{-1}\log2}{1-\kappa}
 \le A+\tau
\]
when $N\ge8/\tau$.  Finally $v_H^{-1}\le8A$ and the definition
of $\kappa$ show that \eqref{eq:critical-channel-width} implies
all three numerical width conditions above.  This proves the theorem.
\end{proof}

\begin{corollary}[Beginning with an approximate finite prefix]
\label{cor:critical-channel-approximate-prefix}
Suppose a controlled finite prefix has relative entropy at most $NE$
against the corresponding original Gaussian blocks.  With probability
at least $1-\kappa/4$, suppose its feature kernel has a measurable
coefficient $c(x)=a^Th(x)$, $\|a\|^2=1/N$, and residual diagonal
$s(x)^2$, such that
\[
 \max_x\sqrt{(c(x)-\sqrt{v_0}f(x))^2+
                            (s(x)-\sqrt{r_0})^2}
 \le\frac{\kappa\sqrt{v_M}}{64e^S}.
\]
Here a scalar-channel continuation has $M\ge1$ remaining hidden
stages, starts at $v_0,r_0$, ends at $v_M$, has score bound $B$, and
has costs $D_1,\ldots,D_M$.  Put
\[
 A=E+\sum_{j=1}^M D_j+\frac1{8v_M},\qquad
 S=\sum_{j=1}^M\sqrt{D_j},\qquad
 \kappa=\min\{1/4,\tau/[32(A+1)]\}.
\]
For widths satisfying \eqref{eq:critical-channel-width} with $H$
replaced by $M$, after a universal change of its constant, the full
network has raw target probability at least $e^{-N(A+\tau)}$.
The prefix and its ideal comparison parameters may depend on $N$,
provided these displayed bounds hold uniformly.
\end{corollary}

\begin{proof}
On a bad prefix use the original subsequent prior.  On a good prefix
run the controlled scalar continuation.  Its metric recursion gives
\[
 e_M\le e^S\{e_0+8M\sqrt\alpha\}.
\]
Use $\sqrt\alpha=\kappa\sqrt{v_M}/(128Me^S)$ and assign failure
probability $\kappa/4$ to each of the prefix, subsequent moment
estimates, and readout noise.  Then $e_M<\kappa\sqrt{v_M}/8$,
and total target success probability is at least $1-\kappa$.
The joint entropy is at most $N(A+3A\kappa)$ by chain rule.
The proof of Theorem~\ref{thm:critical-finite-channel-mass} now
applies without any other change.
\end{proof}

\begin{remark}[Reducing a tracked component]
\label{rem:critical-channel-reduction}
Write $\psi_r(t)=\mathbb E\tanh(t+\sqrt rG)$ and let
$v_{\max}=\mathbb E\psi_r(\sqrt vZ)^2>0$.  The full conditional
mean component uses the score
\[
 b=\psi_r(\sqrt vZ)/\sqrt{v_{\max}}.
\]
Any smaller prescribed component $0<v'\le v_{\max}$ can instead
be obtained with the unit score
\[
 b=\sqrt{\frac{v'}{v_{\max}}}
       \frac{\psi_r(\sqrt vZ)}{\sqrt{v_{\max}}}
       +\sqrt{1-\frac{v'}{v_{\max}}}\,U,
\]
where $U$ is an independent random sign.  Its coefficient is
$\sqrt{v'}f$, and $|b|\le1+(v')^{-1/2}$.  The auxiliary sign adds
no entropy.  This implements every deliberate reduction of a tracked
component in the scalar steering construction.
\end{remark}

\begin{corollary}[Finite-width target mass for Fourier seeds]
\label{cor:critical-finite-fourier-mass}
There are universal $c,C_0,C_1>0$ with the following property.  Put
\[
 A_F(f)=\sum_T|\widehat f(T)|,\qquad
 R=(d+L+2)A_F(f).
\]
If $L\ge C_0\log R$, then for every $N\ge R^c$ the original
critical prior satisfies
\[
 \mu_N(E_f)\ge
 \exp\left[-N\left\{V_*+C_1\frac{\log^2R}{L}\right\}\right].
\]
In particular the sufficient width is polynomial for every parity,
independently of its degree.  This corollary is a target-probability
lower bound, not by itself a posterior learning theorem.
\end{corollary}

\begin{proof}
Use the controlled Fourier mixture with exact mean $f/Z_f$,
$Z_f\le8(d+1)A_F(f)$, and common diagonal.  Its rare component
has probability $\theta=(D_*L)^{-1}$, where
$D_*=C(d+1)^3\log^2(d+2)$.  Retain its Bernoulli mixture tag and
use $b_1=\boldsymbol1_{\{\mathrm{rare}\}}/\sqrt\theta$.
Then $\mathbb E b_1^2=1$, the coefficient is exactly
$\sqrt\theta f/Z_f$, and $v_1=\theta/Z_f^2$.
Using the same component tags in the reference law bounds the joint
row entropy by $\theta D_*\le1/L$.

The established cooling, holding, prescribed growth, purification,
and scalar continuation form a scalar-channel program.  Use the
conditional-mean scores at each stage and
Remark~\ref{rem:critical-channel-reduction} for any deliberate
reduction of the tracked component.  Its total cost satisfies
\[
 A\le V_*+C\log^2R/L.
\]
All tracked positive components are at least $R^{-C}$: the first
one is inverse polynomial; cooling retains at least a constant times
the cooling scale times that component; the holding value is
$v_0=\delta e^{-C S_0}$ with $S_0=O(\log R)$ and
$\delta\ge c/L$; growth and purification retain at least $v_0$.
The final scalar path is increasing from $p\ge c/L$ and has
$O(1/p)$ steps.  Its stable lift retains a fixed positive fraction
of every scalar state, because purification leaves residual
$O(\delta L^{-4})$.  Thus its tracked components are also inverse
polynomial.  All scores are consequently bounded by $B\le R^C$.

The square-root entropy budget also applies when the numbers $D_j$
are upper bounds on the actual feature entropies.  Indeed its proof
uses only scalar dissipation, Bellman suffix lower bounds, and the
total cost; all three remain valid with these larger nonnegative
costs.  Applying that proof to this program gives
\[
 S\le C\log(L+1)+C\sqrt{L(A-V_*)_+}\le C\log R.
\]
Take $\tau=1/L$ in Theorem~\ref{thm:critical-finite-channel-mass}.
Its sufficient width is at most $R^c$ after enlarging a universal
exponent, and its extra $1/L$ cost is absorbed into the displayed
constant $C_1$.
\end{proof}

The finite-width bound is close to the cost of the explicit scalar-channel
construction.  It does not assert a bound within $1/L$ of the unknown
minimum rate $J_L(f)$ for an arbitrary target.  A collective bound on
competing functions is a separate ingredient in any learning consequence.
